% =============================================
% preamble-book.tex
% 用于 R Markdown Book/Article 模板的 LaTeX 导言区
% 支持多语言切换，用法:
%   中文(默认): 不做任何设置
%   英文:       在 YAML header-includes 中加:
%                 header-includes: \newcommand{\envlang}{en}
% =============================================


\usepackage{amsthm}
\usepackage{tcolorbox}
\tcbuselibrary{skins, breakable}


% ---- 默认中文环境名称 ----
\providecommand{\thmname}{定理}
\providecommand{\propname}{命题}
\providecommand{\corname}{推论}
\providecommand{\lemname}{引理}
\providecommand{\claimname}{断言}
\providecommand{\defname}{定义}
\providecommand{\examplename}{例}
\providecommand{\remname}{注}
\providecommand{\notename}{笔记}
\providecommand{\assumename}{假设}
\providecommand{\probname}{问题}
\providecommand{\solutionname}{解}
\providecommand{\conjname}{猜想}
\providecommand{\axiomname}{公理}
\providecommand{\princname}{定律}

% ---- 英文切换（若 header-includes 中定义了 \envlang） ----
% 放在 \AtBeginDocument 中确保可安全使用包命令
\AtBeginDocument{%
  \ifdefined\envlang
    \renewcommand{\thmname}{Theorem}%
    \renewcommand{\propname}{Proposition}%
    \renewcommand{\corname}{Corollary}%
    \renewcommand{\lemname}{Lemma}%
    \renewcommand{\claimname}{Claim}%
    \renewcommand{\defname}{Definition}%
    \renewcommand{\examplename}{Example}%
    \renewcommand{\remname}{Remark}%
    \renewcommand{\notename}{Note}%
    \renewcommand{\assumename}{Assumption}%
    \renewcommand{\probname}{Problem}%
    \renewcommand{\solutionname}{Solution}%
    \renewcommand{\conjname}{Conjecture}%
    \renewcommand{\axiomname}{Axiom}%
    \renewcommand{\princname}{Principle}%
  \fi
}

% ---- 延迟到 \begin{document}，确保 amsthm/hyperref 已加载 ----
\AtBeginDocument{%
  \newtheorem{theorem}{\thmname}[section]
  \newtheorem{proposition}[theorem]{\propname}
  \newtheorem{corollary}[theorem]{\corname}
  \newtheorem{lemma}[theorem]{\lemname}
  \newtheorem{claim}[theorem]{\claimname}
  \theoremstyle{definition}
  \newtheorem{definition}[theorem]{\defname}
  \newtheorem{example}[theorem]{\examplename}
  \theoremstyle{remark}
  \newtheorem{remark}[theorem]{\remname}
  \newtheorem{note}[theorem]{\notename}
  \newtheorem{assumption}[theorem]{\assumename}
  \newtheorem{problem}[theorem]{\probname}
  \newtheorem*{solution}{\solutionname}
  \newtheorem{conjecture}{\conjname}
  \newtheorem{axiom}{\axiomname}
  \newtheorem{principle}{\princname}
  \hypersetup{
    colorlinks = true,
    linkcolor  = blue!70!black,
    urlcolor   = blue!60!black,
    citecolor  = green!50!black,
  }
}

% ---- cbox 彩色盒子 ----
\newtcolorbox{cbox}[1][]{
  colback   = blue!5!white,
  colframe  = blue!75!black,
  coltitle  = white,
  fonttitle = \bfseries,
  title     = {#1},
  breakable,
  boxrule   = 0.8pt,
  arc       = 2pt,
}

% ---- 代码样式 ----
\AtBeginDocument{%
\lstset{
  basicstyle   = \ttfamily\small,
  keywordstyle = \color{blue}\bfseries,
  commentstyle = \color{gray},
  stringstyle  = \color{red},
  numbers      = left,
  numberstyle  = \tiny\color{gray},
  frame        = single,
  breaklines   = true,
  tabsize      = 4,
}
}
